\section*{Recap}
\addcontentsline{toc}{section}{Recap}

In the previous lecture, we characterized an agentic system in terms of five component roles
and an external caller, assigning explicit responsibilities to each.
\begin{itemize}
  \item The \emph{caller} is external to the agent. It initiates or resumes a run, supplies
  the task, input, and applicable constraints, and receives the result.
  \item The \emph{model} interprets observations and proposes subsequent actions through
  reasoning and planning.
  \item The \emph{runtime} is responsible for control flow, validation, scheduling, and
  termination. It executes the control loop, manages the resource budget, and records the
  terminal event.
  \item \emph{Tools} provide the sole mechanism by which proposed actions produce external
  effects, through narrowly scoped, typed interfaces to external systems.
  \item \emph{State} comprises information that persists beyond individual calls, steps, and
  processes. \emph{Context} comprises the information available to the model during a given call.
  \item The \emph{environment} provides the authoritative sources of external data that the
  agent can observe and modify, and encompasses failures that originate outside the agent.
\end{itemize}
Two principles summarize this decomposition. First, the caller initiates the run, the runtime
executes the control loop, and the model proposes subsequent actions. Second, an action request
does not establish that execution has occurred, and execution alone does not establish a
verified outcome.
